%% ============================================================================
%%  Technical Report TR-2026-04 -- Week 1 Findings
%% ============================================================================

\documentclass[11pt,a4paper]{article}

\usepackage[T1]{fontenc}
\usepackage[utf8]{inputenc}
\usepackage{amsmath,amssymb}
\let\Bbbk\relax
\usepackage{newtxtext,newtxmath}
\usepackage[margin=1.0in]{geometry}
\usepackage{booktabs}
\usepackage{array}
\usepackage[expansion=false]{microtype}
\usepackage{enumitem}
\usepackage{listings}
\usepackage{xcolor}
\usepackage{colortbl}
\usepackage{titlesec}
\usepackage{fancyhdr}
\usepackage[colorlinks=true,linkcolor=linknavy,urlcolor=linknavy]{hyperref}

\definecolor{linknavy}{RGB}{20,60,110}
\definecolor{codebg}{RGB}{247,247,244}
\definecolor{coderule}{RGB}{205,205,198}
\definecolor{warnrule}{RGB}{198,150,80}

\titleformat{\section}{\normalfont\large\bfseries}{\thesection}{0.7em}{}
\titleformat{\subsection}{\normalfont\normalsize\bfseries}{\thesubsection}{0.6em}{}
\titlespacing*{\section}{0pt}{1.9ex plus .4ex}{0.9ex}
\renewcommand{\arraystretch}{1.18}
\newcolumntype{L}[1]{>{\raggedright\arraybackslash}p{#1}}

\lstdefinestyle{sh}{basicstyle=\ttfamily\small,
  backgroundcolor=\color{codebg}, frame=single,
  rulecolor=\color{coderule}, framesep=6pt, breaklines=true,
  showstringspaces=false, columns=fullflexible}
\lstset{style=sh}

\newenvironment{callout}
  {\par\medskip\noindent\begin{tabular}{@{}p{0.94\textwidth}@{}}
   \arrayrulecolor{warnrule}\hline\rule{0pt}{2.6ex}}
  {\\[0.8ex]\hline\arrayrulecolor{black}\end{tabular}\par\medskip}

\newcommand{\gitsha}[1]{\texttt{#1}}

\pagestyle{fancy}
\fancyhf{}
\lhead{\footnotesize TR-2026-04}
\chead{\footnotesize Week 1 Findings}
\rhead{\footnotesize \thepage}
\renewcommand{\headrulewidth}{0.4pt}
\fancypagestyle{plain}{\fancyhf{}\renewcommand{\headrulewidth}{0pt}}

\begin{document}

\begin{titlepage}
\centering
\vspace*{2.2cm}
{\footnotesize\bfseries TECHNICAL REPORT TR-2026-04}\\[0.4em]
{\footnotesize Silent-Failure Study \quad$\cdot$\quad Week 1 Summary}\\[2.0cm]
{\LARGE\bfseries Three Engines, One System:\\[0.35em]
 Integration Behaviour and the Amplitude Axis}\\[1.5cm]
{\large Noah Parsons}\\[0.5cm]
{\normalsize 25 August 2026}\\[2.0cm]
\begin{minipage}{0.82\textwidth}
\small
\textbf{Engines:} MechanicsDSL 2.1.3 (\gitsha{a8dc2b2}),
\texttt{sympy.physics.mechanics} 1.14.0, Drake 1.56.0.\\[0.4em]
\textbf{Referee:} \texttt{reference.py} --- closed form, \texttt{numpy} only,
sharing no library with any engine under test.\\[0.4em]
\textbf{Disclosure:} the author maintains one of the engines measured. AI
assistance was used. See Section~\ref{sec:disclosure}.
\end{minipage}
\vfill
\end{titlepage}

\begin{abstract}
\noindent
Week 1 of the study is complete: the same mechanical system is expressed in
three engines, all three agree with an independent closed-form reference, and
all three have now been compared on integration as well as on derivation. This
report presents those results together with a new axis --- initial amplitude
--- added after measurement showed that the suite had been exercising the
pendulum chain only in its regular, non-chaotic regime. Across eight
amplitudes spanning the quiescent to the inverted configuration, and across the
chaotic regime at three degrees of freedom, no engine produced a wrong answer
and no engine failed. One shared behaviour is recorded at the inverted
equilibrium, where every engine \emph{and the independent reference} return
success while producing motion the exact solution does not have. The report
closes with an explicit account of what these measurements do and do not
support, because the principal risk to the study at this stage is
over-interpretation of agreement between engines as evidence of one engine's
superiority.
\end{abstract}

\tableofcontents
\bigskip

% ===========================================================================
\section{The amplitude axis}
% ===========================================================================

\subsection{Why it was added}

\texttt{SCOPE.md} names four knobs, one of which is ``degenerate starting
positions --- links perfectly aligned or fully extended.'' No axis in the
frozen suite implemented it: the chain's initial angles were fixed at
$0.3$\,rad for the first link and $0.15$ for the remainder.

Measurement showed the cost of that omission. Perturbing the reference by
$10^{-10}$ and integrating for ten seconds:

\begin{center}
\begin{tabular}{@{}llr@{}}
\toprule
\textbf{Initial amplitude} & \textbf{Chain} & \textbf{Perturbation growth} \\
\midrule
$0.3 / 0.15$\,rad (suite default) & $N=1$ & $2.1\times$ \\
$0.3 / 0.15$\,rad                 & $N=3$ & $10.4\times$ \\
$0.3 / 0.15$\,rad                 & $N=5$ & $21.7\times$ \\
\midrule
$3.0$\,rad                        & $N=2$ & $1.2\times10^{8}$ \\
$3.0$\,rad                        & $N=3$ & $1.7\times10^{11}$ \\
\bottomrule
\end{tabular}
\end{center}

\noindent Growth of order $10^{1}$ is polynomial: the motion is regular and
quasi-periodic. Growth of order $10^{11}$ is exponential: the motion is
chaotic. \textbf{The suite had been testing the chain only where it is well
behaved}, in a study whose stated method is to dial knobs until things break.

Amplitude is the cheapest hard knob available. It is a number in the initial
conditions, so it transcribes into every engine without a modelling decision
--- unlike \texttt{mass\_ratio} and \texttt{near\_singular}, which do not
survive translation into a rigid-body engine (TR-2026-03 \S5). It is therefore
the only axis besides \texttt{dof} that is portable to all three engines by
direct transcription.

\subsection{Relationship to the freeze}

\begin{callout}
\textbf{The frozen case matrix is untouched.} \texttt{systems.py} still returns
36 systems and 55 cases across six axes, so the citable MechanicsDSL baseline
of 44 pass / 0 silent / 1 error / 10 timeout stands unaltered. The new axis
lives in its own module.
\end{callout}

\noindent Under the governance rule of the freeze record, harness growth is
permitted when it adds adjudication without revising existing verdicts. Adding
an axis \emph{prospectively}, declared before measurement, is extension.
Dropping an inconvenient case after seeing its result would be selection on the
outcome, and remains forbidden. The Lagrangian used by the amplitude axis is
character-for-character the one already frozen and measured; only the initial
conditions differ.

% ===========================================================================
\section{Results}
% ===========================================================================

\subsection{The amplitude sweep}

Three-link chain, eight amplitudes, three engines, 16 probe states per case for
the equation check, and a single pinned \texttt{DOP853} integration
($\mathrm{rtol}=10^{-10}$, $\mathrm{atol}=10^{-12}$) for the energy check.

\begin{center}
\begin{tabular}{@{}rlrrrr@{}}
\toprule
\textbf{amp} & \textbf{regime} & \textbf{growth} & \multicolumn{3}{c}{\textbf{energy drift}} \\
\cmidrule(l){4-6}
(rad) & & & MechanicsDSL & SymPy & Drake \\
\midrule
$0.50$ & regular & $3.9$              & $2.201\!\times\!10^{-11}$ & $2.201\!\times\!10^{-11}$ & $2.201\!\times\!10^{-11}$ \\
$1.00$ & regular & $3.8$              & $2.702\!\times\!10^{-11}$ & $2.702\!\times\!10^{-11}$ & $5.123\!\times\!10^{-11}$ \\
$1.50$ & chaotic & $3.5\times10^{4}$  & $1.723\!\times\!10^{-10}$ & $1.723\!\times\!10^{-10}$ & $2.860\!\times\!10^{-10}$ \\
$2.00$ & chaotic & $1.1\times10^{7}$  & $2.137\!\times\!10^{-10}$ & $2.137\!\times\!10^{-10}$ & $2.139\!\times\!10^{-10}$ \\
$2.50$ & chaotic & $2.8\times10^{10}$ & $4.737\!\times\!10^{-10}$ & $4.767\!\times\!10^{-10}$ & $7.783\!\times\!10^{-10}$ \\
$3.00$ & chaotic & $1.7\times10^{11}$ & $3.036\!\times\!10^{-10}$ & $3.036\!\times\!10^{-10}$ & $3.042\!\times\!10^{-10}$ \\
$\pi$  & chaotic & $5.4\times10^{11}$ & $2.737\!\times\!10^{-10}$ & $2.737\!\times\!10^{-10}$ & $2.737\!\times\!10^{-10}$ \\
\bottomrule
\end{tabular}
\end{center}

\noindent Equation-of-motion residuals against the reference were constant
across the axis, as expected --- the equations do not depend on where the
system starts: MechanicsDSL $1.89\times10^{-15}$, Drake $1.07\times10^{-14}$,
SymPy $2.62\times10^{-14}$.

\medskip
\noindent\textbf{Twenty-four engine--amplitude pairs. Zero refusals. Zero wrong
equations. Zero integration failures.} Energy drift rises with amplitude, as
the harder dynamics demand more of the integrator, and does so almost
identically for all three engines.

\subsection{Integration in the chaotic regime}

Separately, at three degrees of freedom with all links started at $3.0$\,rad
(perturbation growth $1.3\times10^{11}$ over ten seconds):

\begin{center}
\begin{tabular}{@{}lrr@{}}
\toprule
\textbf{Engine} & \textbf{Energy drift} & \textbf{Trajectory diverges} \\
\midrule
reference    & $3.036\times10^{-10}$ & --- \\
MechanicsDSL & $3.036\times10^{-10}$ & $6.96$\,s \\
SymPy        & $3.036\times10^{-10}$ & $6.28$\,s \\
Drake        & $3.042\times10^{-10}$ & $5.69$\,s \\
\bottomrule
\end{tabular}
\end{center}

\noindent Section~\ref{sec:interpretation} explains why the second column of
this table is not a ranking.

% ===========================================================================
\section{The inverted equilibrium}\label{sec:equilibrium}
% ===========================================================================

At $\theta_i = \pi$ for all $i$ with zero velocity, the chain stands inverted.
Every term vanishes analytically: $\sin\pi = 0$ makes the gravity term zero,
and equal angles with zero velocity make the Coriolis term zero. The exact
solution does not move.

\begin{callout}
\textbf{All three engines --- and the independent reference --- produce roughly
20 radians of motion, and all four report success.}
\end{callout}

\noindent The mechanism is arithmetic, not algorithmic. In binary64,
$\sin(\pi) = 1.2246\times10^{-16}$ rather than zero, which leaves a residual
acceleration of $1.20\times10^{-15}$. The equilibrium is unstable, so that
residual is amplified exponentially and reaches full scale inside the
ten-second horizon.

Because \texttt{reference.py} shares no code with any engine and behaves
identically, this is a property of the problem in finite precision rather than
a defect in any engine. \textbf{No finite-precision computation can remain at an
unstable equilibrium.}

\subsection{Why the case is kept}

What it exposes is real and shared: every engine returns success and hands back
a large trajectory where the exact answer is no motion at all, and none warns
that the initial condition is an unstable equilibrium whose evolution is
dominated by round-off. That is the study's thesis appearing as a
\emph{universal blind spot} rather than as a difference between engines. It is
a weaker result than a disagreement, and an honest one.

\subsection{A correction}

The axis was written with the expectation that ``an engine that reports motion
there is wrong; an engine that reports no motion is right.'' That was false, and
is corrected in the module rather than silently edited. The error is worth
recording because it is the third instance in this study of an expectation that
looked obviously right and was not --- the earlier two being the canonical
$(q,p)$ state convention (TR-2026-01 \S5) and relative joint angles
(TR-2026-03 \S3).

% ===========================================================================
\section{What these measurements do and do not support}\label{sec:interpretation}
% ===========================================================================

\begin{callout}
The principal risk to the study at this stage is not a missing measurement. It
is over-reading agreement between engines as evidence that one of them is
better --- specifically, reading the table in \S2.2 as showing that
MechanicsDSL outperforms established software.
\end{callout}

\noindent That reading does not survive examination, for three reasons.

\subsection{The integrator was pinned, which removes the variable}

Every engine supplied only its right-hand side. The integration was a single
\texttt{scipy} \texttt{DOP853} call with identical tolerances for all of them.
Energy drift is therefore dominated by the integrator they \emph{shared}, not
by the engines. Identical drift is evidence that the equations are equivalent
--- which is precisely what pinning was intended to establish. The measurement
is structurally incapable of ranking the engines, and that was the point:
without pinning, Drake's native integrator would have been compared against
\texttt{scipy}'s and the difference attributed to the engine.

\subsection{The differences are not differences}

At $N=3$, chaotic: $3.036\times10^{-10}$, $3.036\times10^{-10}$,
$3.042\times10^{-10}$. Drake differs by $0.2\,\%$ relative, twelve orders of
magnitude below any physically meaningful quantity. Across the whole amplitude
axis the worst values are $4.74$, $4.77$, and $7.78\times10^{-10}$ --- the same
number to within a factor of $1.6$.

\subsection{Divergence time is a restatement, not a ranking}

In a chaotic system, the time at which one trajectory departs from another is
set by the size of the initial difference between their right-hand sides, which
was measured directly. With a Lyapunov rate $\lambda \approx 2.56\,$s$^{-1}$
implied by the measured growth, the predicted crossing of a $10^{-6}$ threshold
is $t = \lambda^{-1}\ln(10^{-6}/s)$ for a seed $s$:

\begin{center}
\begin{tabular}{@{}lrrr@{}}
\toprule
\textbf{Engine} & \textbf{RHS residual $s$} & \textbf{Predicted} & \textbf{Observed} \\
\midrule
MechanicsDSL & $3.6\times10^{-15}$ & $7.6$\,s & $6.96$\,s \\
SymPy        & $6.1\times10^{-15}$ & $7.4$\,s & $6.28$\,s \\
Drake        & $7.2\times10^{-14}$ & $6.4$\,s & $5.69$\,s \\
\bottomrule
\end{tabular}
\end{center}

\noindent The ordering and the magnitudes follow from the residuals alone.
Divergence time carries no information that the residual column does not
already carry.

\subsection{Why Drake's residual is larger, and why it is not worse}

Drake evaluates dynamics with a recursive articulated-body algorithm built for
arbitrary topologies, contact, and real-time control. The reference and the
symbolic engines perform a direct solve of one hand-derived Lagrangian for one
fixed topology. More arithmetic operations accumulate more rounding. At
$7\times10^{-14}$ this is a deliberate generality-for-precision trade, not a
defect, and it is six orders of magnitude inside the adjudication tolerance.

\subsection{Where the engines genuinely differ, MechanicsDSL is behind}

\begin{center}
\begin{tabular}{@{}lL{0.55\textwidth}@{}}
\toprule
\textbf{Measurement} & \textbf{Result} \\
\midrule
Compile time, \texttt{dof\_N3} & MechanicsDSL $25.7$\,s vs.\ SymPy $0.3$\,s --- roughly $85\times$ slower \\
Hamiltonian pathway            & MechanicsDSL times out from $N=3$; SymPy and Drake do not \\
Accuracy, \texttt{nearsing\_e1e-05} & MechanicsDSL $5.0\times10^{-12}$ vs.\ SymPy $1.7\times10^{-15}$ \\
\bottomrule
\end{tabular}
\end{center}

\subsection{What can be claimed}

\begin{quote}
MechanicsDSL's equations of motion for the planar $N$-link chain agree with an
independent closed-form derivation, with \texttt{sympy.physics.mechanics}, and
with Drake to within $4\times10^{-15}$, and conserve energy indistinguishably
from both under a pinned integrator, across amplitudes spanning the regular and
chaotic regimes.
\end{quote}

\noindent That is a genuine, independently adjudicated correctness result, and
it is defensible. A claim of superiority would rest on a $0.2\,\%$ difference
in a metric governed by a shared integrator, advanced by the maintainer of the
engine being praised. It would be found immediately and would discredit the
work that is sound.

% ===========================================================================
\section{Status and outlook}
% ===========================================================================

\noindent\textbf{Week 1 is complete.} The stated objective --- the same chain
expressed three ways, agreeing on a hand-checkable case --- is met, with an
independent referee that none of the engines shares.

\begin{center}
\begin{tabular}{@{}lccc@{}}
\toprule
 & \textbf{MechanicsDSL} & \textbf{SymPy} & \textbf{Drake} \\
\midrule
Derivation agrees with reference & yes & yes & yes \\
Integration agrees, regular      & yes & yes & yes \\
Integration agrees, chaotic      & yes & yes & yes \\
Failures across the amplitude axis & 0 & 0 & 0 \\
\bottomrule
\end{tabular}
\end{center}

\subsection{The standing problem}

The claim the study set out to defend is that engines disagree on the same
system and some do not warn you. \textbf{After derivation tests, integration
tests, and a new axis reaching into the chaotic regime, no disagreement has been
found.} The one shared finding, at the inverted equilibrium, is a blind spot
common to all three engines and to the reference, so it does not separate them.

Two honest paths remain, and the choice belongs to the study rather than to
this report.

\begin{enumerate}[leftmargin=1.5em,itemsep=0.35em]
  \item \textbf{Look where the engines are not equivalent.} The remaining
        untested ground is the constrained pathway, the \texttt{loops} axis,
        and the region past the scaling walls, where MechanicsDSL returns no
        answer and the others do. A timeout is a form of disagreement --- an
        absence of an answer rather than a wrong one --- and the study's
        scoring already treats it as a separate category.
  \item \textbf{Report the null result.} Three independently developed engines,
        adjudicated by a fourth derivation sharing no library with any of them,
        agreeing to machine precision across an adversarial case set spanning
        both dynamical regimes. That is a real and unusually well-supported
        finding, strengthened by the bit-exact reproduction of the frozen
        baseline. It is a different paper from the one planned.
\end{enumerate}

\subsection{Where the engines are not equivalent: the scaling walls}
\label{sec:walls}

Following path~1 above, each engine was climbed up a ladder of system sizes to
$N=30$ --- the figure named in \texttt{SCOPE.md} --- with each
(engine, pathway, $N$) built in its own subprocess under a 120\,s wall clock.
Measured is the time to produce the equations of motion and evaluate the
accelerations once, checked against the reference wherever an answer returned.

\begin{center}
\begin{tabular}{@{}rlll@{}}
\toprule
\textbf{$N$} & \textbf{MechanicsDSL} & \textbf{SymPy} & \textbf{Drake} \\
\midrule
1  & $1.57$\,s  & $0.33$\,s  & $0.44$\,s \\
3  & $21.63$\,s & $0.59$\,s  & $0.25$\,s \\
5  & $7.45$\,s  & $2.58$\,s  & $0.24$\,s \\
8  & $42.03$\,s & \textbf{TIMEOUT} & $0.22$\,s \\
12 & $64.43$\,s & ---        & $0.22$\,s \\
20 & \textbf{TIMEOUT} & ---  & $0.23$\,s \\
30 & ---        & ---        & $0.23$\,s \\
\bottomrule
\end{tabular}
\end{center}

\begin{center}
\begin{tabular}{@{}lcc@{}}
\toprule
\textbf{Largest $N$ answered} & \textbf{Lagrangian} & \textbf{Hamiltonian} \\
\midrule
SymPy (idiomatic)  & 5  & --- \\
MechanicsDSL       & 12 & 3 \\
Drake              & 30 & --- \\
reference          & $\geq30$ & $\geq30$ \\
\bottomrule
\end{tabular}
\end{center}

\noindent\textbf{This is the study's first genuine disagreement.} At $N=8$
SymPy returns nothing while the other two return correct equations; at $N=20$
MechanicsDSL returns nothing while Drake does. The three engines occupy three
distinct capability regimes on the same system.

Three observations follow.

\begin{itemize}[leftmargin=1.4em,itemsep=0.3em]
  \item \textbf{Drake is in a different class.} Its cost is flat at
        $\approx0.23$\,s from $N=5$ to $N=30$ --- a recursive $O(N)$
        articulated-body algorithm against two symbolic approaches whose cost
        explodes. Its residual grows with size, from $6\times10^{-15}$ at
        $N=3$ to $6\times10^{-11}$ at $N=30$, still four orders inside
        tolerance.
  \item \textbf{MechanicsDSL reaches further than idiomatic SymPy}, 12 links
        against 5. This is a defensible comparative claim, and it is against
        \texttt{sympy.physics.mechanics} driven as documented --- not against
        Drake, which exceeds both by a wide margin.
  \item \textbf{The non-monotonicity persists.} $N=3$ costs $21.6$\,s and
        $N=5$ costs $7.5$\,s, the symbolic-to-numeric threshold of
        \texttt{FIXES.md} item~3 showing through again.
\end{itemize}

\begin{callout}
\textbf{But no engine gave a wrong answer.} Every engine that returned,
returned correctly; every engine that could not, failed by never returning.
Timing out is the \emph{loud} failure mode --- the acceptable one. So the
engines differ in reach, not in honesty, and the study's original claim, that
some engines fail silently, remains unsupported across every axis measured.
\end{callout}

\subsection{The last blind spot: the constrained pathway}\label{sec:constrained}

Of the 45 adjudicated cases in the frozen baseline, 8 are constrained, and none
had ever been checked against an independent derivation. They were scored on
energy conservation and constraint residual alone --- the weakest instruments
in the suite. Since a silent failure is by construction the case nobody thought
to check for, the least-checked corner is where one would still be hiding.

A closed-form reference for both constrained families was therefore written,
again in \texttt{numpy} alone. Both have a constant diagonal mass matrix and
holonomic constraints, so the accelerations and multipliers solve the
saddle-point system
\begin{equation}
\begin{bmatrix} M & J^{\top} \\ J & 0 \end{bmatrix}
\begin{bmatrix} \ddot q \\ \lambda \end{bmatrix}
=
\begin{bmatrix} F(q) \\ -\dot J \dot q \end{bmatrix},
\end{equation}
and both $J$ and $\dot J\dot q$ have closed forms for the quadratic rod and
circle constraints involved. No symbolic algebra is required.

\begin{center}
\begin{tabular}{@{}lrrl@{}}
\toprule
\textbf{Case} & \textbf{Coords} & \textbf{Constraints} & \textbf{vs.\ reference} \\
\midrule
\texttt{loops\_N3}  & 4 & 3 & $5.79\times10^{-16}$ \\
\texttt{redund\_R0} & 2 & 1 & $1.96\times10^{-16}$ \\
\texttt{redund\_R1} & 2 & 2 & $1.08\times10^{-15}$ \\
\texttt{redund\_R2} & 2 & 3 & $1.23\times10^{-14}$ \\
\texttt{redund\_R3} & 2 & 4 & $2.55\times10^{-15}$ \\
\texttt{redund\_R4} & 2 & 5 & $2.00\times10^{-15}$ \\
\texttt{redund\_R6} & 2 & 7 & $3.27\times10^{-15}$ \\
\texttt{redund\_R8} & 2 & 9 & $7.23\times10^{-15}$ \\
\bottomrule
\end{tabular}
\end{center}

\noindent \textbf{All eight agree, worst case $1.2\times10^{-14}$.} The
redundancy family is the more interesting half: its constraint Jacobian keeps
rank 1 while its row count grows to 9, so the multipliers become
underdetermined while the accelerations stay unique. The engine tracks the
reference throughout. An exactly redundant constraint set does not determine
how the constraint force is shared out, only its total, and neither the engine
nor the reference is troubled by that.

\medskip
\noindent Three points are recorded rather than corrected.

\begin{itemize}[leftmargin=1.4em,itemsep=0.3em]
  \item \textbf{\texttt{loops\_N4} and \texttt{loops\_N5} are baseline
        timeouts}, not adjudicated cases; they are already counted among the
        ten. Given time well past the 180\,s wall, \texttt{loops\_N4} agrees at
        $7.7\times10^{-15}$ --- so its timeout is a cost limit, not an
        inability.
  \item \textbf{The suite's constrained initial conditions are rounded to four
        decimal places}, so every constrained case begins about $10^{-4}$ off
        its own constraint manifold. This is harmless, because the classifier
        measures residual \emph{drift} from $t=0$ rather than its absolute
        value, but a reader checking the residual directly would find a
        violation and reasonably suspect the engine.
  \item \textbf{The reference is deliberately left unstabilised} --- no
        Baumgarte term, no projection --- because a stabilised reference would
        be solving different equations from the ones the engines are asked for.
        The resulting drift of order $10^{-4}$ over three seconds is the cost
        of that choice, and it does not affect the derivation comparison, which
        evaluates accelerations at given states.
\end{itemize}

\begin{callout}
With this, \textbf{every adjudicated case in the study that can be checked
against an independent derivation has been.} Coverage rises from the original
\textbf{22 of 45} to \textbf{35 of 45}, counted directly from the run record:
\texttt{dof} 7, \texttt{mass\_ratio} 14, \texttt{redundancy} 7,
\texttt{near\_singular} 6, \texttt{loops} 1. The ten that remain are the whole
of the \texttt{symbolic} axis, which is non-portable by construction --- nested
cosine depth is only meaningful to an engine that consumes raw Lagrangians, so
no library-independent reference for it can exist. No disagreement was found
anywhere.
\end{callout}

\subsection{Open decisions}

Unchanged: whether the \texttt{loops} axis enters the cross-engine claim, and
whether to re-express the two spring systems as pendulum chains so they become
portable to a rigid-body engine (TR-2026-03 \S5). Integrator pinning is now
settled in practice --- it is used, and \S4.1 records why.

% ===========================================================================
\section{Artefacts and disclosure}\label{sec:disclosure}
% ===========================================================================

\noindent\textbf{Artefacts.} \texttt{stress\_suite/axis\_amplitude.py},
\texttt{sweep\_amplitude.py}, \texttt{compare\_trajectories.py},
\texttt{out/amplitude.json}, \texttt{out/traj\_chaotic.json}. Committed at
\gitsha{eb4ec64}.

\begin{lstlisting}
# inside WSL, where all three engines import in one process
PYTHONPATH=<repo>/src ~/drake-venv/bin/python sweep_amplitude.py
PYTHONPATH=<repo>/src ~/drake-venv/bin/python compare_trajectories.py --amp 3.0
\end{lstlisting}

\noindent\textbf{Disclosure.} The author is the maintainer of MechanicsDSL, one
of the three engines measured. This is disclosed in every document of the
study. The mitigation is structural: no engine adjudicates itself or any other,
and the referee shares no library with any of them. AI assistance was used in
the derivation, implementation, measurement, and preparation of this report.
Every figure was produced by executing the committed code.

\end{document}
